\scanpage{102}

le \unsure{schéma|résultat} correct pour le $G/G^\circ$: car en effet, il est dual au groupe des racines de l'unité, or ce dernier est bien le dual de $\widehat Z$!

\margin{\[
M(q^m)\subset M(\underset{\shortparallel\,p^n}{q})\subset M(p)
\]
\[
M(q^m)^m\subset M(q)^m\subset M(q^m)\subset M(q)
\]}

Soit $M(q)$ le s-gr. de $Q'(i)^*$ qu'on définit en termes de $q$ ci-dessus, soit \deleted{$\unsure{T}$} $\inserted{f_q\in}G(q)$ son dual, qui est le groupe de Galois motivique sur $\mathbb F_q$. Considérons l'extension $\mathbb F_{q^m}$ de $\mathbb F_q$, d'où
\deleted{$\unsure{G_q^m}$}
\diagram{\begin{tikzcd}[column sep=large]
G(q)&G(q^m)\arrow[l]\\
(f_q)^m&f_{q^m}\arrow[l,rightsquigarrow]
\end{tikzcd}}

\margin{$M(q)\simeq Z\times M^\circ\deleted{(q)}$ si $q=p^{\unsure{n}}$, \unsure{$n$ pair}. $M^\circ\deleted{(q)}$ formé des $u\in Q'(i)^*$ qui sont des unités abs. exc. $p$, et vérifient $v_{\mathfrak p}(u)+v_{\mathfrak p'}(u)=0$; le premier facteur \deleted{\unsure{$Z$}} $Z$ correspond à $q^{1/2}$. \unsure{Les homomorphismes de transition entre les} $M(q)$ \unsure{sont} l'identité sur le facteur $Z$, et la multiplication par $m$ sur le facteur $M^\circ$.}

\deleted{\unsure{de fait}} on voit que l'hom désigné \unsure{$\alpha$} est induit par
\diagram{\begin{tikzcd}[column sep=huge]
M(q)\arrow[r,"{\text{élévation à la puissance }m}"]&M(q^m)
\end{tikzcd}}
dont le noyau est comme de juste \deleted{$\unsure{Z}$} $\mu_m$ (dual de $Z/mZ$) et dont l'image est $M(q)^m$ donc contient $M(q^m)^m$, dont le conoyau est annulé par $m$. Par suite, le \deleted{conoyau} noyau de $\alpha$ est annulé par $m$.

\deleted{\unsure{Comme}} Considérons la limite inductive des $M(q)$ ($q$ variant multiplicativemt) par les hom précédents: on trouve $M(p^\infty)$, un groupe sans torsion dont le dual doit être le groupe de

\mspage{101}
